\documentclass[conference]{support/IEEEtran}
\IEEEoverridecommandlockouts

% IEEE conference version of the ACT Workshop draft (same text and figure
% placeholders). Body sections live in sections/, references in references/,
% and the IEEEtran class/bibliography style in support/.
\usepackage{cite}
\usepackage{amsmath,amssymb,amsfonts}
\usepackage{graphicx}
\usepackage{textcomp}
\usepackage[table]{xcolor}
\usepackage{booktabs}
\usepackage{multirow}
\usepackage{tikz}
\usetikzlibrary{arrows.meta,positioning,calc,fit}
\usepackage{algorithm}
\usepackage{algpseudocode}
\usepackage[hidelinks]{hyperref}

\newcommand{\system}{\textsc{AS-Depth}}
\definecolor{ACMRed}{HTML}{BF0000}
\newcommand{\todo}[1]{\textcolor{ACMRed}{[TODO: #1]}}
\definecolor{paperblue}{HTML}{24577A}
\definecolor{paperlight}{HTML}{EEF3F7}
\definecolor{papergreen}{HTML}{347264}
% acmart's accessibility description has no IEEEtran counterpart.
\providecommand{\Description}[1]{}
\newcommand{\pending}{\textcolor{paperblue}{\textit{Pending}}}
\newcommand{\draftnote}[1]{\par\smallskip\noindent
  {\small\color{paperblue}\textit{Draft note: #1}}\par\smallskip}
\newcommand{\figslot}[3]{%
  \IfFileExists{#1}{\includegraphics[width=\linewidth,height=#2,keepaspectratio]{#1}}{%
  \fcolorbox{paperblue!45}{paperlight}{%
  \parbox[c][#2][c]{\dimexpr\linewidth-2\fboxsep-2\fboxrule\relax}{%
  \centering\color{paperblue}\small #3\par\medskip
  {\footnotesize Figure reserved for the advisor-review draft}}}}}
\newcommand{\resultslot}[2]{\fcolorbox{paperblue!45}{paperlight}{%
  \parbox[c][#1][c]{\dimexpr\linewidth-2\fboxsep-2\fboxrule\relax}{%
  \centering\small\color{paperblue}#2\par\medskip
  \footnotesize Measurements pending; no synthetic results shown.}}}
% Placeholder for a table whose rows and columns are not fixed yet.
\newcommand{\tabslot}[2]{\fcolorbox{paperblue!45}{paperlight}{%
  \parbox[c][#1][c]{\dimexpr\linewidth-2\fboxsep-2\fboxrule\relax}{%
  \centering\small\color{paperblue}#2\par\medskip
  \footnotesize Table layout pending; no synthetic results shown.}}}

\begin{document}

\title{AS-Depth: Unleashing OFA Collaboration For Efficient Metric Depth Completion on SoCs}

\author{\IEEEauthorblockN{Anonymous Author(s)}
\IEEEauthorblockA{Anonymous Institution\\
Anonymous}}

\maketitle

% Original draft (kept for reference):
% Recent breakthrough in Depth completion models enables the reconstruction of dense metric depth from sparse or noisy measurements captured by low-cost sensors,
% offering great benefits to autonomous driving and robotics.
% Existing methods achieve promising accuracy with large Transformer-based architectures, which can make them diffcult to deploy for real-time near-sensor perception due to theTransformer-heavy computation and large parameter counts.
% However, industrial robotic systems often require (1)real-time inference (2)deployment on resource-constrained platforms (e.g., SoC).
% To address these limitations, we analyse the scheduling model of current depth completion pipeline and observe large temporal redundancies of continuous depth maps under high framerates.
% Building on this insight, we propose as-depth, an asynchronous framework dedicated for accelerating depth completion on SoCs via three techniques: (1) an asynchronous path which decouples the frame display and model computation stages (2) a heteregeous hardware collaboration featured in Optical Flow Accelerator on SoC (3) an adaptive frame adjusment strategy.
% Extensive experiments on depth completion tasks demonstrate over 6 $\times$ speedup and 50\% energy savings. We further show that as-depth can be integrated into otherexisting depth completion frameworks.
\begin{abstract}
Depth completion recovers dense metric depth from the sparse and noisy
measurements of low-cost RGB-D sensors, and has become a key enabler for
robotic manipulation and autonomous driving.
Recent Transformer-based models have achieved promising accuracy on metric depth prediction, 
but their heavy computation and large parameter counts confine them to
run real-time framerate on edge platforms.
While modern robotic systems demand real-time depth on resource-constrained
edge SoCs.
To address this limitation, we observe that this gap lies in the model itself and how the
pipeline is scheduled: every frame is pushed through the full network, even
though consecutive depth maps at high frame rates are highly redundant.
Building on this insight, we present \system{}, an asynchronous framework to accelerate depth completion on SoCs through three techniques:
(1)~an asynchronous execution path that decouples depth output from model
inference;
(2)~a heterogeneous collaboration that offloads inter-frame depth propagation
to the SoC's on-chip Optical Flow Accelerator (OFA) and
(3)~an adaptive frame-adjustment strategy that tunes the inference rate at
runtime.
On NVIDIA Jetson AGX Thor, \system{} delivers over
6$\times$ speedup and 50\% energy savings with negligible accuracy drop. 
Moreover, \system{} can integrate into existing depth-completion
models.
\end{abstract}

\begin{IEEEkeywords}
Computer systems organization → Sensors and actuators; Computing methodologies → Perception; Vision for robotics.
\end{IEEEkeywords}

\section{Introduction}
\label{sec:introduction}

Metric depth measures the distance of each pixel in an image from the camera, 
providing geometry cues such as spatial and shape information that the 2D color 
image cannot provide. 
Such depth information can be inputted into the backbone model of 
applications such as autonomous driving and robotics shown in Fig.~\ref{fig:teaser},
which greatly improves
the capability of spatial reasoning and enables precision interaction
with the objects.

With the widespread availability of commodity depth cameras, acquiring depth information appears to be low cost and straightforward. 
However, commodity RGB-D cameras frequently suffer from missing values, 
depth holes, flying pixels, and measurement noise, 
especially under challenging materials such as transparent, 
reflective, and specular surfaces. 
Although recent studies have integrated depth into robotic manipulations, performance remain limited by the poor quality of depth data produced by such device. Therefore, evaluations of robotics is typically restricted to simulators where clean depth is available. 
Thus, depth completion emerges as an effective method for generating dense metric depth from commodity depth cameras, and has become a critical component for bridging the Sim to Real gap.

% Teaser: its source sits in text set in the right column of page 1, so the
% [t] float lands on top of that column instead of above the abstract.
\begin{figure}[t]
  \centering
  \figslot{figures/teaser.pdf}{1.30in}{%
    Teaser\\[5pt]
    Raw RGB-D with missing/noisy depth $\rightarrow$ completed metric depth}
  \caption{Teaser. \pending}
  \label{fig:teaser}
\end{figure}

% Full pipeline (Fig. 2): a figure* is set at the top of the page after the
% one holding its source, so it stays on page 1 to land on top of page 2.
% It must follow the teaser: LaTeX keeps figure and figure* in source order,
% so an earlier deferred figure* would push the teaser to page 2.
\begin{figure*}[t]
  \centering
  \figslot{figures/pipeline.pdf}{2.00in}{%
    Complete AS-Depth pipeline\\[5pt]
    RGB-D input $\rightarrow$ ViT encoder / optical flow $\rightarrow$
    heterogeneous scheduling on the SoC $\rightarrow$ completed metric depth}
  \caption{Overview of the complete \system{} pipeline. \pending}
  \label{fig:pipeline}
\end{figure*}

Recent studies 1234 have demonstrated advanced accuracy of depth completion, but their increasingly large Transformer-based architectures impose substantial computation costs. The selected GPU baselines range from 4.26 to 6.62 FPS on RTX 3090/4090 platforms. Because the models and platforms differ, these figures provide context rather than a controlled speedup comparison. Achieving real-time depth completion on an edge SoC such as NVIDIA Jetson AGX Thor remains challenging under its computation budget, motivating algorithm and hardware acceleration.

% To mitigate the contradiction of model heavy computation cost and SoC's computation budget, we propose a novel method: offloading one part of the existing model's computation to another device --- specifically, a FPGA --- to accelerate the whole depth completion pipeline. By leveraging the FPGA's flexible hardware acceleration capabilities, we aim to achieve real-time performance on edge platforms. We begin by profiling the depth completion pipeline and the operator type, and then analyzing the computational workload distribution between the SoC and the FPGA to find the optimal partitioning strategy. We also explore the best customized-precision strategies for the FPGA implementation, and maintain the model accuracy throughout quantization aware training. 

To address this, we propose As-Depth, a asynchronous heterogeneous acceleration system built on SoC and FPGA tailored for depth completion task on edge platforms, as shown in Fig.~\ref{fig:pipeline}.
As-Depth achieves efficient real time rendering speed and SOTA rendering quality on the edge platform. At the architecure level, it decouples the image representation encoding from the metric depth decoding stage, enabling these two tasks to be executed asynchronously on the SoC and FPGA according to their computational characteristics. On the algorithmn-hardware co-design level, our depth aware quantization and customized decoder design effectively preserve the accuracy of metric depth while reducing computational overhead. At the accelerator level, our depth aware accelerator integrates an ethernet interface unit and a depth decoding engine to handle the decoding pipeline, enabling efficient data communication and computation.

Building on these insights above, we develop a novel heterogeneous acceleration system, As-Depth, that combines the strengths of both the SoC and the FPGA to achieve real-time depth completion on edge platforms. It leverages the mature CUDA cores and Tensor Cores of the SoC to accelerate the transformer-based encoder, while offloading the CNN decoder to our customized FPGA implementation. This design ensures the maximum utilization of both the SoC and the FPGA, enableing pipelined execution to hide the latency of the offloaded operations. Through our quantization aware training and depth aware customized precision, our system strike a balance between performance and resource overhead. Our contributions are as follows:  

\begin{itemize}
  \item We propose a novel heterogenous system and combine task level asynchronous execution for depth completion based on the existing depth completion pipeline, efficiently mapping computing units according to the different computation characteristics of the SoC and FPGA, and using asynchronous optical queued communication to improve inference performance. 
  \item We design a depth-aware quantization strategy, and a customized precision scheme for the fpga implementation. Our approach successfully balance the accuracy of metric depth between model inference latency.
  \item We implement a prototype of our As-Depth system on the Nvidia Jetson AGX Thor and Intel Stratrix FPGA platform. Our prototype runs at real-time speed of 30FPS, demonstrating the feasibility of our heterogeneous acceleration approach.
  \item We implement and evaluate As-Depth on Synopsys. Experiments show that our system achieves 5.4x speedup and xxx improvement in energy efficiency over the latest SOTA implementation on the targeted SoC with the whole depth completion pipeline. To the best of our knowledge, this is the first work that achieves real-time 30FPS inference of depth completion on edge platforms.
  Our system enables real-time depth near sensor possible which contribute to mitigate the Sim to Real gap in the future rbotics applications. 
  Moreover, our approach demonstrates the potential for heterogeneous acceleration in resource-constrained environments, paving the way for future advancements in edge AI systems.
\end{itemize}

\section{Related Works}
\label{sec:related}

\subsection{Depth Completion}
\label{sec:related-depth}

Depth completion estimates a dense metric-depth map from an RGB image and an
incomplete depth observation~\cite{uhrig2017sparsity}.  The input depth may
contain holes, outliers, and material-dependent noise, whereas the output must
retain both absolute scale and local geometry for downstream robotic
perception.  Modern methods commonly pair a representation encoder with a
dense prediction decoder.  The encoder extracts global visual and geometric
features; the decoder fuses and upsamples these features to the target depth
resolution.

\subsection{ViT Compression}
\label{sec:related-vit}

Prior GPU depth-completion systems establish the accuracy of large models, and
prior FPGA accelerators demonstrate the efficiency of specialized low-precision
datapaths.  \system{} connects these directions by retaining the encoder on the
SoC, specializing the decoder on FPGA, and treating the model boundary and
cross-frame schedule as first-class parts of the design.
\draftnote{General ViT compression techniques, e.g.,
\cite{jacob2018quantization}; contrast with \system{}, whose optimizations are
application specific to depth completion.}

\section{Method}
\label{sec:method}

\system{} combines three mechanisms derived from the observations below.  A
task-level execution framework partitions the model and overlaps consecutive
frames; DAQ coordinates low precision across the Thor encoder and FPGA decoder;
and a dedicated decoder accelerator turns the regular decoder workload into an
efficient hardware dataflow.

\subsection{Preliminary}
\label{sec:preliminary}

% Observation figure: top of the left column on page 3.
\begin{figure}[t]
  \centering
  \figslot{figures/observation.pdf}{1.30in}{%
    Observation figure}
  \caption{Observation. \pending}
  \label{fig:obs}
\end{figure}

\system{} preserves the encoder--decoder decomposition of depth completion and
uses its boundary as a system partition.  The selected depth-completion network
and its trained weights are the workload being accelerated, while this paper
focuses on the low-precision deployment, heterogeneous execution, and dedicated
decoder hardware.

\paragraph{Computational characteristics}
The two stages expose different opportunities for edge acceleration.  The
encoder contains Transformer blocks whose matrix operations and MLP layers can
use Thor's mature GPU and Tensor Core execution stack.  The decoder is composed
primarily of structured dense-prediction operators for which customized
parallelism, buffering, and precision can be implemented on FPGA.  This
difference motivates a stage-level mapping rather than repeatedly moving
individual operators between devices.  We profile the unmodified model on Thor
and report the latency of preprocessing, encoder, decoder, and postprocessing
under a fixed resolution, precision, batch size, and power mode
(Fig.~\ref{fig:obs} and Table~\ref{tab:obs}).

\paragraph{Edge SoC--FPGA heterogeneous execution}
An edge SoC offers a complete software environment and high-throughput tensor
units, but its compute and energy budget is shared with sensing, planning, and
control.  An FPGA provides configurable precision, dataflow, and memory
organization for a stable subgraph.  The combination is useful only when the
partition is coarse enough to limit communication and when execution across
devices is scheduled to hide, rather than merely add, transfer latency.

\paragraph{Observation 1: Monolithic execution misses the real-time budget}
The Transformer encoder and CNN decoder have different operator mixes and
hardware preferences, yet a conventional deployment executes both stages
sequentially on one processor.  The resulting end-to-end latency includes the
full encoder, decoder, and runtime path for every frame.  Merely moving the
decoder to another device does not remove this serialization if communication
and accelerator execution remain on the critical path.  The clear
encoder--decoder boundary instead permits a coarse-grained partition: the Thor
executes the encoder while the FPGA executes the decoder.  Across a stream of
frames, queued communication and double buffering can overlap the encoder of
frame $i+1$ with transfer and decoding of frame $i$.

% Observation table: top of the right column on page 3.
\begin{table}[t]
  \caption{Observation table. \pending}
  \label{tab:obs}
  \centering
  \tabslot{0.80in}{Observation table}
\end{table}

\paragraph{Observation 2: End-to-end low precision can degrade metric depth}
Low precision is beneficial on both sides of the partition.  In particular,
selected MLP computation in the Thor-side encoder and the FPGA-side decoder
can be executed at reduced precision.  A uniform quantization policy, however,
must simultaneously accommodate Transformer activation distributions,
depth-regression sensitivity, and the numerical contract at the cross-device
boundary.  Aggressive quantization can therefore save computation and
bandwidth while introducing scale errors or local depth artifacts.  Rather than
treating encoder quantization, feature transfer, and decoder quantization as
independent optimizations, \system{} coordinates their precision choices.

\subsection{Our Design: Optical Flow}
\label{sec:optical-flow}

\draftnote{Optical-flow design, illustrated in Fig.~\ref{fig:optical-flow}.}

% Optical-flow figure: bottom of the right column on page 3.
\begin{figure}[b]
  \centering
  \figslot{figures/optical-flow.pdf}{1.30in}{%
    Optical-flow design}
  \caption{Optical-flow design. \pending}
  \label{fig:optical-flow}
\end{figure}

\paragraph{Depth-aware quantization}
DAQ treats quantization as an end-to-end property of the heterogeneous model.
On Thor, it targets selected MLP computation in the encoder; at the partition
boundary, it defines the representation transmitted to the FPGA; and inside
the FPGA, it determines the precision of decoder weights, activations,
accumulators, and output depth.  Training optimizes the quantized path against
metric-depth supervision so that numerical error is weighted by its effect on
the final depth map rather than by feature reconstruction alone, extending the
standard quantization-aware training setting~\cite{jacob2018quantization}.

\paragraph{Dedicated depth decoder accelerator}
The FPGA accelerator specializes the weighted and auxiliary operators in the
depth decoder while retaining a programmable layer schedule.  Its datapath
must expose the precision contract established by DAQ, reuse weights and
activations on chip where profitable, and stream larger intermediate tensors
through external memory without changing model semantics.  The communication
front end converts received feature packets into accelerator-visible buffers
and returns the completed depth map after the final layer.  This design follows
the broader principle of mapping quantized networks to customized FPGA
dataflows~\cite{umuroglu2017finn} while specializing the execution path for
metric-depth decoding.  For every layer, the scheduler records tensor
dimensions, addresses, quantization parameters, and dependencies, and
verification compares each hardware layer boundary against the quantized
software reference.

\subsection{Scheduling in Heterogeneous System}
\label{sec:scheduling}

\paragraph{Model partition and tensor contract}
The model is partitioned once at the encoder--decoder boundary.  Thor consumes
the RGB image and sensor depth, executes the encoder, and exports a quantized
feature tensor.  The FPGA consumes this tensor and returns the completed
metric-depth map.

\paragraph{Queued cross-frame execution}
The host runtime assigns separate buffers to encoder output, outbound
transfer, decoder input, and completed depth.  Producer--consumer queues make
ownership explicit and prevent a buffer from being reused before its consumer
finishes.  After pipeline fill, the encoder for the next frame can progress in
parallel with transfer and decoder execution for the current frame, as
summarized in Algorithm~\ref{alg:scheduling}.

% Scheduling pseudocode: bottom of the left column on page 4.
\begin{algorithm}[b]
  \caption{Heterogeneous scheduling (placeholder)}
  \label{alg:scheduling}
  \begin{algorithmic}[1]
    \Require Input frame stream \pending
    \Ensure Completed metric-depth stream \pending
    \For{each incoming frame}
      \State \pending
    \EndFor
  \end{algorithmic}
\end{algorithm}

\paragraph{Communication and back pressure}
The runtime batches feature metadata with the boundary tensor and preserves
frame order on the return path.  If either device becomes the bottleneck,
bounded queues apply back pressure instead of silently dropping or reordering
frames.  Section~\ref{sec:experiments} measures transfer time separately from
useful computation so that communication hiding can be verified.

\section{Experiments}
\label{sec:experiments}

We evaluate \system{} by answering four questions.  First, can heterogeneous
execution sustain the 30-FPS real-time target once communication and runtime
overheads are included?  Second, does DAQ preserve metric-depth quality on
challenging sensor data?  Third, how does the complete system compare with
recent depth-completion methods?  Finally, what resource, communication,
power, and energy costs are required to obtain these gains?

\subsection{Settings}
\label{sec:settings}

\paragraph{Baselines}
The controlled system comparison runs the same depth-completion network in
three deployments: \emph{GPU-only} on an RTX~4090, \emph{Thor-only} on the
target edge SoC, and \system{}.  The GPU-only and Thor-only baselines execute
the complete model sequentially on one device; \system{} includes feature
transfer and output reconstruction in its end-to-end time.  For model-level
context, we additionally compare against CDM, PriorDA, PromptDA, and OMNI-DC,
using their released checkpoints and official evaluation procedures whenever
available.  Because their published FPS values come from different GPU
platforms, we label the execution platform of every method and do not treat
cross-platform FPS as a controlled speedup claim.

\paragraph{Benchmarks}
We evaluate all accuracy-oriented experiments on HAMMER and ClearPose.  These
datasets exercise different object geometries and difficult sensing conditions,
so results are reported for each dataset rather than averaged into one score.
Every method uses the same test split, validity mask, crop, input resolution,
and depth range within a dataset.  The system and ablation experiments use
batch size one, which matches continuous RGB-D capture on the target robotic
platform.

\paragraph{Experiment platform}
The prototype comprises an NVIDIA Jetson AGX Thor and an Intel Stratix FPGA
connected through the system interface described in
Section~\ref{sec:scheduling}.  The FPGA clock, external-memory configuration,
link rate, Thor power mode, and software versions are held fixed across all
\system{} variants.  The RTX~4090 baseline uses the same model weights, input
resolution, preprocessing, and output validation.

\paragraph{Metrics and measurement protocol}
We measure both single-frame latency and steady-state throughput.  Latency is
decomposed into preprocessing, encoder execution, boundary-tensor transfer,
decoding, return transfer, and output reconstruction; FPS is measured only
after the asynchronous pipeline is filled.  Timing includes synchronization
needed for a valid completed-depth output, rather than reporting kernel time
alone.  Depth quality is evaluated with AbsRel, RMSE, MAE, and
$\delta_{1.05}$.  Lower is better for the three error metrics and higher is
better for $\delta_{1.05}$.  We report system power, energy per completed
frame, and FPS/W under a consistent measurement boundary.

\subsection{Main Experiments}
\label{sec:main-exp}

% Main results (table*): its source stays on page 4 so it lands on top of
% page 5.
\begin{table*}[t]
  \caption{Main results. Values marked ``--'' await measurement.}
  \label{tab:main-exp}
  \centering
  \small
  \setlength{\tabcolsep}{4pt}
  \renewcommand{\arraystretch}{1.10}
  \begin{tabular}{llc|ccc|ccc}
    \toprule
    & & & \multicolumn{3}{c|}{HAMMER} & \multicolumn{3}{c}{ClearPose} \\
    Method & Platform & FPS $\uparrow$
      & AbsRel $\downarrow$ & RMSE $\downarrow$ & $\delta_{1.05}$ $\uparrow$
      & AbsRel $\downarrow$ & RMSE $\downarrow$ & $\delta_{1.05}$ $\uparrow$ \\
    \midrule
    CDM & -- & -- & -- & -- & -- & -- & -- & -- \\
    PriorDA & -- & -- & -- & -- & -- & -- & -- & -- \\
    PromptDA & -- & -- & -- & -- & -- & -- & -- & -- \\
    OMNI-DC & -- & -- & -- & -- & -- & -- & -- & -- \\
    \system{} & -- & -- & -- & -- & -- & -- & -- & -- \\
    \bottomrule
  \end{tabular}
\end{table*}

\paragraph{Throughput}
Table~\ref{tab:main-exp} compares the inference rate of CDM, PriorDA,
PromptDA, OMNI-DC, and \system{}, together with the execution platform of
each method.  This comparison answers a deployment question---whether current
depth-completion methods can service a real-time sensor stream---while the
controlled GPU-only/Thor-only comparison supports architectural speedup
claims.  We report \system{}'s measured end-to-end FPS and identify whether its
advantage comes from the coarse model partition, cross-frame overlap, or
low-precision execution, using the ablation evidence rather than the
cross-platform numbers alone.

\paragraph{Depth accuracy}
Table~\ref{tab:main-exp} also compares all methods under common HAMMER and
ClearPose protocols.  The discussion reports the absolute metric for every
method, \system{}'s gap to the best-performing method, and its rank on each
dataset.  The relevant question is whether \system{} reaches real-time edge
throughput while keeping the accuracy loss small and consistent across both
datasets.

\subsection{Ablation Studies}
\label{sec:ablation}

\paragraph{Optical flow versus frequency}
\draftnote{Ablation on optical flow at different rates, shown in
Fig.~\ref{fig:of-vs-hz}.}

% OF vs. Hz figure: bottom of the left column on page 5.
\begin{figure}[b]
  \centering
  \resultslot{1.30in}{Optical flow vs.\ Hz}
  \caption{Ablation on optical flow versus frequency. \pending}
  \label{fig:of-vs-hz}
\end{figure}

\paragraph{Scheduling}
\draftnote{Ablation on scheduling, reported in Table~\ref{tab:abl-scheduling}.}

\paragraph{Heterogeneous execution}
Thor-only versus \system{} is the primary controlled comparison for the
heterogeneous partition: both process the same frames with the same network,
while only the latter offloads the decoder and overlaps encoder,
communication, and decoder work across frames.  We therefore report not only
their FPS difference, but also the stage-latency breakdown and the fraction of
transfer hidden by the asynchronous schedule.  GPU-only is retained as a
high-end performance reference, not as an edge-device energy baseline.

\paragraph{Metric-depth accuracy}
The decisive accuracy comparison is \system{} without versus with DAQ because
their model partition, target bit widths, and runtime are identical.  Any
quality difference therefore measures the value of coordinating the quantized
encoder MLP, boundary representation, and decoder using metric-depth
supervision.  We report each dataset independently and discuss whether
improvements come from reducing large depth errors, increasing the fraction of
accurate pixels, or both.

% Temporary: the draft body does not fill the left column of page 5 yet, so
% end it here; Tables III and IV are then read in the right column. Remove
% once the ablation text fills the left column.
\newpage

% Scheduling ablation: right column on page 5, below Table II.
\begin{table}[t]
  \caption{Ablation on scheduling. \pending}
  \label{tab:abl-scheduling}
  \centering
  \tabslot{0.80in}{Ablation on scheduling}
\end{table}

% Fallback table: bottom of the right column on page 5.
\begin{table}[b]
  \caption{Fallback optical flow on CUDA. \pending}
  \label{tab:fallback}
  \centering
  \tabslot{0.80in}{Fallback optical flow on CUDA}
\end{table}

\paragraph{Fallback optical flow on CUDA}
\draftnote{Optical flow falls back to CUDA; results in Table~\ref{tab:fallback}.}

\paragraph{Power and energy efficiency}
Power is measured for Thor-only and the complete heterogeneous system using
the same input stream and steady-state window used for FPS.  We derive energy
per frame as system power divided by completed frames per second and
additionally report FPS/W.  The analysis separates the extra power introduced
by the additional hardware from the energy saved by shortening the end-to-end
execution time.

% Visualization (figure*): its source stays on page 5 so it lands on top of
% page 6.
\begin{figure*}[t]
  \centering
  \figslot{figures/visualization.pdf}{2.40in}{%
    Qualitative visualization with zoom-in regions}
  \caption{Visualization. \pending}
  \label{fig:visualization}
\end{figure*}

% Temporary: the draft body does not fill page 5 yet, so start page 6 here to
% keep Fig. 6 above Detailed Visualization. Remove once page 5 is full.
\newpage

\subsection{Detailed Visualization}
\label{sec:visualization}

\draftnote{Detailed visualization in Fig.~\ref{fig:visualization}, followed by
a zoom-in explanation.}

\section{Conclusion}
\label{sec:conclusion}

This paper presents \system, a heterogeneous depth-completion system that
decouples representation encoding on Thor from depth decoding on FPGA.  Its
task-level runtime overlaps consecutive frames, while DAQ coordinates
low-precision encoder MLP computation, the cross-device representation, and
the dedicated decoder datapath.  The final manuscript will close by reporting
the measured throughput, metric-depth accuracy, and energy efficiency achieved
by this combination on HAMMER and ClearPose.


\bibliographystyle{support/IEEEtran}
\bibliography{references/as-depth,references/architecture}
\end{document}
